\documentclass[10pt]{article}
\usepackage{titling}

\title{Equilibrium}
\date{DATE PLACEHOLDER}
\author{Ingyu Bahng}

\usepackage[letterpaper, margin=1in]{geometry}
\usepackage{amsmath} % better math functions
\usepackage{amssymb} % better math symbols
\usepackage{babel}[english] % change rules to english
\usepackage{lipsum} % allows for lorem ipsum text generation
\usepackage{xr-hyper} % allows cross referencing labels from external LaTeX documents 
\usepackage{hyperref} % for hyperlinking text
\usepackage{fancyvrb} % for better verbatim environments
\usepackage{newverbs} % allows for inline typewriter font via \texttt{}
\usepackage{fancyvrb} % also code font blocks but has more capabilities than texttt
\usepackage{footnote} % allows footnotes inside tabulars
\usepackage{dcolumn} % allows for decimal alignment in tables
\usepackage{tabularx}

\usepackage{fontspec}
\setmainfont{Gentium Plus} % allowing greek letters in text

\usepackage{unicode-math}
\setmathfont{XITS Math} % allowing greek letters in math



% tcolorbox package
\usepackage{tcolorbox}
\tcbuselibrary{breakable}

\tcbset{ % this is the default design of tcolorboxes
  colframe = black,
  colback  = black!1,
  coltitle = black,
  colbacktitle = black!15,
  breakable, 
  arc=0mm,
  boxrule=0.5pt,
  toptitle=3pt,
  bottomtitle=3pt,
  left=8pt,
  right=8pt,
  top=6pt,
  bottom=6pt,
  before skip=12pt,
  after skip=12pt,
  bottomrule at break=-1pt,
  toprule at break=-1pt,
  fonttitle=\bfseries,
}

% Custom Environments
\newtcolorbox[auto counter, number within=section]{definition}[1][]
{
    title = \textbf{Definition \thetcbcounter: ~#1}
}

\newtcolorbox[auto counter, number within=section]{core}[1][]
{
    title = \textbf{Core Concept \thetcbcounter: ~#1}
}


\usepackage{fancyhdr}
\usepackage[skip=0.8em, indent=0em]{parskip}

\pagestyle{fancy} % this section generates the lines and text at the top and bottom of each page
\fancyhead[L]{\thetitle}
\fancyhead[C]{\theauthor}
\fancyhead[R]{\thedate}
\fancyfoot[C]{\thepage}
\renewcommand{\footrulewidth}{0.3pt}
\renewcommand{\thispagestyle}[1]{}

\renewcommand{\arraystretch}{1.5} % table vertical padding
\setlength{\tabcolsep}{10pt} % table horizontal padding

\usepackage{enumitem} % better control over enumerate and itemize

% List customization
\setlist[itemize]{%
  label=\bullet,
  itemsep=2pt,
  leftmargin=1.5em,
}

\setlist[enumerate]{%
  label=(\alph*),
  itemsep=2pt,
  leftmargin=1.5em,
}



\usepackage{pgfplots} % for plot creation
\pgfplotsset{compat=1.18} % setting the version used for pgfplots
\usepackage{float} % for better figure placement using [h]
\usepackage{graphicx} % for importing figures into the tex file
\usepackage{subcaption} % allows for multiple subfigures with individual captions wihtin one figure
\usepackage{xparse} % allows for more than 9 parameters in commands
\usepackage{adjustbox} % for valign option
\captionsetup{font=small} % setting the caption fontsize to small always

\usepackage{silence}
\WarningsOff[float]  % suppress float package warnings

\newcommand{\myfig}[4]{%
  \begin{figure}[H]
    \centering
    \adjustbox{valign=c}{\includegraphics[width=#1\textwidth]{#2}}
    \caption{#3}
    \label{#4}
  \end{figure}
}

\newcommand{\mymultifig}[8]{
  \begin{figure}[H]%
    \centering%
    \begin{subfigure}{#1\textwidth}%
      \centering%
      \includegraphics[width=\linewidth]{#2}%
    \end{subfigure}%
    \hfill%
    \begin{subfigure}{#3\textwidth}%
      \centering%
      \includegraphics[width=\linewidth]{#4}%
    \end{subfigure}%
    \hfill%
    \begin{subfigure}{#5\textwidth}%
      \centering%
      \includegraphics[width=\linewidth]{#6}%
    \end{subfigure}%
    \caption{#7}%
    \label{#8}%
  \end{figure}%
}

\usepackage{newverbs} % allows for inline typewriter font via \texttt{}
\usepackage{fancyvrb} % also code font blocks but has more capabilities than texttt
\usepackage{listings} % allows for codeblocks - config below

%BELOW ARE CODE BLOCK DESIGNS (Light Theme)
\definecolor{gray}{HTML}{707070}
\definecolor{lightgray}{HTML}{333333}
\definecolor{codebg}{HTML}{F5F5F5}
\definecolor{keyword}{HTML}{0000FF}
\definecolor{string}{HTML}{A31515}
\definecolor{number}{HTML}{098658}
\definecolor{function}{HTML}{795E26}
\definecolor{comment}{HTML}{008000}
\definecolor{classname}{HTML}{267F99}

% default options for listings (for code)
\lstset{
  frame=ltbr,
  language=python,
  aboveskip=3mm,
  belowskip=3mm,
  showstringspaces=false,
  columns=fullflexible,
  keepspaces=true,
  basicstyle=\fontsize{9}{10}\ttfamily\color{lightgray},
  numbers=left,
  firstnumber=1,
  numberstyle=\fontsize{7}{10}\ttfamily\color{gray},
  keywordstyle=\color{keyword},
  commentstyle=\color{comment},
  stringstyle=\color{string},
  backgroundcolor=\color{codebg},
  breaklines=true,
  breakatwhitespace=true,
  tabsize=2,
  xleftmargin=3em,
  numbersep=1em,
  framexleftmargin=2.5em,
  framesep=6pt,
  stepnumber=1,
}

\hfuzz=5.002pt % ignore overfull hbox badness warnings below this limit



\usepackage{chemfig}
\usepackage[version=4]{mhchem}

\newcommand{\bce}[1]{
  \hspace{7mm}{#1}
}

\newcommand{\sno}[0]{
  $\mathrm{S}_{\mathrm{N}}1$
}

\newcommand{\snt}[0]{
  $\mathrm{S}_{\mathrm{N}}2$
}

\newcommand{\reactionfig}[4]{
    \begin{figure}[H]
    \centering
    \scalebox{#4}{
    \schemestart
    #1
    \schemestop
    }
    \caption{#2}
    \label{#3}
    \end{figure}
}




\externaldocument[gas-]{../gases/gases}[../gases/gases.pdf]
\externaldocument[aci-]{../acids_and_bases/acids_and_bases}[../acids_and_bases/acids_and_bases.pdf]
\externaldocument[the-]{../thermodynamics/thermodynamics}[../thermodynamics/thermodynamics.pdf]
\externaldocument[qua-]{../quantization/quantization}[../quantization/quantization.pdf]
\externaldocument[alc-]{../alcohols/alcohols}[../alcohols/alcohols.pdf]
\externaldocument[eth-]{../ethers/ethers}[../ethers/ethers.pdf]
\externaldocument[nom-]{../nomenclature/nomenclature}[../nomenclature/nomenclature.pdf]
\externaldocument[rad-]{../radical_halogenation/radical_halogenation}[../radical_halogenation/radical_halogenation.pdf]
\externaldocument[con-]{../conformational_analysis/conformational_analysis}[../conformational_analysis/conformational_analysis.pdf]
\externaldocument[alk-]{../alkane_alkene_alkyne/alkane_alkene_alkyne}[../alkane_alkene_alkyne/alkane_alkene_alkyne.pdf]
\externaldocument[nuc-]{../nuc_sub_and_elim/nuc_sub_and_elim}[../nuc_sub_and_elim/nuc_sub_and_elim.pdf]
%\externaldocument[equ-]{../equilibrium/equilibrium}[../equilibrium/equilibrium.pdf]

\begin{document}

\input{../../presets/_general/startup}

\section{Reaction Rate \& Equilibrium}

  \begin{core}[label=placeholder1]{Reaction Rates}

    Atoms, ions, and molecules react with one another via collisions. The more frequent and energetic the collisions, the faster the reaction rate.

    Take the formation of compound \ce{C} as an example: \ce{A + B -> C}. The rate at which \ce{C} forms depends on factors such as[A], [B], temperature, the presence of a catalyst, etc. and can be expressed mathematically as: 
    \[
      \text{Rate} = k[\ce{A}]^m[\ce{B}]^n
    \]

    Here, k is the \textbf{rate constant}, which is specific to the given reaction and its conditions. \textit{m} and \textit{n} are the \textbf{partial reaction orders}, which represents a reactant's influence over the reaction rate. Here is a hypothetical plot of [A], [B], and [C] over time:

    \pgfig{
      \begin{axis}
        [
          width=9cm,
          height=7cm,
          xlabel={Time (s)},
          ylabel={Concentration (mol dm$^{-3}$)},
          xmin=0, xmax=10,
          ymin=0, ymax=4.5,
          xtick=\empty,
          ytick={1,2,3,4},
          xlabel style={at={(axis description cs:0.5,-0.08)}},
          ylabel style={at={(axis description cs:-0.12,0.5)}},
          \axispresets
        ]

        \addplot[
          domain=0:10,
          samples=200,
          thick,
          red,
        ]{2.8 + 1.2*exp(-0.4*x)};
        \node[red] at (axis cs:10.3, 2.8) {A};

        \addplot[
          domain=0:10,
          samples=200,
          thick,
          green!70!black,
        ]{2.3 + 1.2*exp(-0.4*x)};
        \node[green!70!black] at (axis cs:10.3, 2.3) {B};

        \addplot[
          domain=0:10,
          samples=200,
          thick,
          violet,
        ]{1.6 - 1.3*exp(-0.4*x)};
        \node[violet] at (axis cs:10.3, 1.6) {C};

      \end{axis}
    }{Equilibrium of \ce{A + B -> c}}{filler1}{0.8}
    
    In single-step reactions, partial reaction orders are equal to the stoichiometric coefficients of in the balanced chemical equation. However, in multi-step reactions, partial reaction orders may not directly correspond to the stoichiometric coefficients and must be determined experimentally.

  \end{core}
  
  Consider the decomposition of \ce{H2O2} ($\ce{H2O2 -> H2O + O2}$) as an example. When you graph the concentration of \ce{H2O2} over time, you get an exponential decay curve, where the demposition rate of \ce{H2O2} decreases as [\ce{H2O2}] decreases.

  \pgfig{
    \begin{axis}
      [
        width=9cm,
        height=7cm,
        xlabel={Time (h)},
        ylabel={[H$_2$O$_2$] (mol L$^{-1}$)},
        xmin=0, xmax=26,
        ymin=0, ymax=1.1,
        xtick={0,6,12,18,24},
        ytick={0.200,0.400,0.600,0.800,1.000},
        yticklabel style={/pgf/number format/fixed},
        \axispresets
      ]

      \addplot[
        thick,
        black,
        domain=0:25,
        samples=300,
      ]{exp(-0.12*x)};

    \end{axis}
  }{Decomposition of \ce{H2O2} over time}{filler2}{0.8}

  \begin{core}[label=placeholder2]{Equilibrium}
    
    Reactions do not always proceed in only one direction. Many reactions are \textbf{reversible}, meaning that the products can react to reform the reactants. When the rates of the forward and reverse reactions are equal, the system is said to be at \textbf{equilibrium}. Thus, for the general reversible shell reaction \ce{AB <=> A + B}, the Eq reaction state can be expressed as:
    \[
      k_{forward}[\ce{AB}] = k_{backward}[\ce{A}][\ce{B}]
    \]

    And the equilibrium constant $K_c$ is defined as:
    \[
      K_c = \frac{k_f}{k_b} = \frac{[\ce{A}][\ce{B}]}{[\ce{AB}]}
    \]

  \end{core}

  \begin{core}{Analyzing $K_c$}
    
    $K_c$ provides insight into the position of equilibrium. A large $K_c$ ($>>1$) indicates product favorability, while a small $K_c$ ($<<1$) indicates reactant favorability.

    Some examples of reactions that possess a \textbf{large $K_c$} (reaction to completion) include:

    \bce{\ce{Si + O2 -> SiO2}\hspace{5mm}$K_c \approx 10^{150}$}

    \bce{\ce{2H2 + O2 -> H2O}\hspace{5mm}$K_c \approx 10^{81}$}

    Some examples of reactions that possess a \textbf{moderate $K_c$} (half-reaction) include:

    \bce{\ce{CO + Hemoglobin -> CO*Hemoglobin}\hspace{5mm}$K_c \approx 300$}

    \bce{\ce{O2 + Hemoglobin -> O2*Hemoglobin}\hspace{5mm}$K_c \approx 4$}

    \bce{\ce{N2O4 -> 2NO2}\hspace{5mm}$K_c \approx 0.2$}

    An example of a reaction that possess a \textbf{neglible $K_c$} (no reaction) include:

    \bce{\ce{N2 + O2 -> 2NO}\hspace{5mm}$K_c \approx 10^{-31}$}

  \end{core}

  The decomposition of Bromothymol Blue (\textit{denoted as \ce{A}}) exhibits a nice visual representation of reaction equilibria.
 
  \bce{\ce{\textcolor{Yellow}{HA} <=> H+ + \textcolor{Blue}{A-}}}

  If we start with a pure HA solution, the solution will gradually change color from \textcolor{Yellow}{yellow} to \textcolor{Green}{green} to \textcolor{Blue}{blue} as the forward reaction proceeds and [\ce{A^-}] rises. This particular compound is also known as an indicator, a concept covered in \autoref{aci-def:indicator}.
  
  \begin{definition}[label=placeholder3]{Le Chatelier's Principle}

    \textbf{Le Chatelier's Principle} states that a system in equilibrium subjected to a stress will react in a direction that tends to reduce that stress and return back to Eq.

    In more exact terms, stress will shift the \textbf{reaction quotient (Q)}---which represents the relative amounts of reactant and products at a given point in time before it has nessecarily reached equilibrium ($Q \neq K$)---and the system will react in a manner that shifts Q back to K.
    \[
      Q = \frac{[\ce{A}]_{current}[\ce{B}]_{current}}{[\ce{AB}]_{current}}
    \] 

  \end{definition}

  There are three major categories of stress in the context of Le Chatelier's Principle.

  \begin{core}{Le Chatelier's Principle (\textit{Concentration Change})}
   
    Shifting the point concentrations of any chemical participant will shift Q.

    Let's reexamine the Bromothymol Blue reaction (\ce{\textcolor{Yellow}{HA} <=> H+ + \textcolor{Blue}{A-}}). If we add more \ce{H+} to an initially \textcolor{Green}{green} solution ( $[\textcolor{Yellow}{\ce{HA}}] = [\textcolor{Blue}{\ce{A-}}]$ ), the system will increase the backwards reaction rate to consume excess \ce{H+}, resulting in a more \textcolor{Yellow}{yellow} solution. Conversely, removing \ce{H+} (\textit{by adding \ce{OH-}}) will increase the forward reaction rate to  produce more \ce{H+}, resulting in a more \textcolor{Blue}{blue} solution.

    Consider another example (\ce{PCl5 -> PCl3 + Cl2}) based on the visual below. The plot shows the abrupt addition of \ce{PCl5} to a system originally at equilibrium. Consequently, $Q<K$ and the system reacts by increasing the forwards reaction rate to re-establish equilibrium, resulting in a subsequent increase in [\ce{PCl3}] and [\ce{Cl2}].

    \pgfig{
      \begin{axis}
        [
          width=11cm,
          height=7cm,
          xlabel={time},
          ylabel={concentration},
          xmin=0, xmax=10,
          ymin=0, ymax=5.5,
          xtick=\empty,
          ytick=\empty,
          x axis line style={-stealth},
          y axis line style={-stealth},
          \axispresets
        ]

      % --- left equilibrium zone (0 to 3.5) ---

      % PCl5 — top flat line
      \addplot[thick, blue, domain=0:3.5, samples=2]{3.0};

      % PCl3 — middle flat line
      \addplot[thick, black, domain=0:3.5, samples=2]{2.5};

      % Cl2 — bottom flat line
      \addplot[thick, red!70!black, domain=0:3.5, samples=2]{1.5};

      % --- disturbance zone (3.5 to 6.5) ---

      % PCl5 — drops then flattens
      \addplot[thick, blue, domain=3.5:6.5, samples=100]
        {5.0 - 0.8*(1 - exp(-1.5*(x-3.5)))};

      % PCl3 — rises then flattens
      \addplot[thick, black, domain=3.5:6.5, samples=100]
        {2.5 + 0.6*(1 - exp(-1.5*(x-3.5)))};

      % Cl2 — rises then flattens
      \addplot[thick, red!70!black, domain=3.5:6.5, samples=100]
        {1.5 + 0.6*(1 - exp(-1.5*(x-3.5)))};

      % --- right equilibrium zone (6.5 to 10) ---

      % PCl5
      \addplot[thick, blue, domain=6.5:10, samples=2]{4.2};

      % PCl3
      \addplot[thick, black, domain=6.5:10, samples=2]{3.1};

      % Cl2
      \addplot[thick, red!70!black, domain=6.5:10, samples=2]{2.1};

      % --- labels on right ---
      \node[blue, font=\small]          at (axis cs:10.3, 4.2) {[PCl$_5$]};
      \node[black, font=\small]         at (axis cs:10.3, 3.1) {[PCl$_3$]};
      \node[red!70!black, font=\small]  at (axis cs:10.3, 2.1) {[Cl$_2$]};

      % --- Q labels at bottom ---
      \node[font=\small] at (axis cs:1.75, -1.0)  {$Q = K_C$};
      \node[font=\small] at (axis cs:5.0,  -1.0)  {$Q < K_C$};
      \node[font=\small] at (axis cs:8.25, -1.0)  {$Q = K_C$};

      % --- disturbance vertical dashed line ---
      \draw[dashed, gray] (axis cs:3.5, 0) -- (axis cs:3.5, 5.5);
      \draw[dashed, gray] (axis cs:6.5, 0) -- (axis cs:6.5, 5.5);

      % --- annotation arrow ---
      \draw[-stealth, red!80]
        (axis cs:4.5, 5.4) -- (axis cs:3.6, 5.1);
      \node[red!80, font=\footnotesize, align=center]
        at (axis cs:6.3, 5.4)
        {equilibrium disturbed by\\changing conc. of PCl$_5$};

      \end{axis}
    }{Disturbance of equilibrium by changing concentration of \ce{PCl5}}{fig:pcl5equilibrium}{1}

  \end{core}

  \begin{core}{Le Chatelier's Principle (\textit{Pressure \& Volume Change})}

    While concentration is the primary focus for solutions, pressure and volume are the key variables for gaseous systems. Increasing the pressure forces the equilibrium toward the side with fewer gas molecules to relieve the stress. Conversely, decreasing pressure allows the reaction to shift toward the side with more gas molecules.

    The equilibrium constant in terms of partial pressures for gaseous reactions is written as \ce{K_p}.

    Consider the following reaction:

    \bce{\ce{2NO2(g) <=> N2O4(g)}\hspace{5mm}$K_p = \frac{P_{N2O2}}{(P_{NO2})^2} \approx 7$ at 25°C}

    Increasing the volume ($V$) causes a drop in total pressure ($P$). To counteract this, the equilibrium shifts to the left, favoring the side with more moles of gas. Mathematically, this occurs because the partial pressure of \ce{NO2} is squared in the denominator of the $Q_p$ expression. A decrease in pressure impacts the squared term more significantly than the linear numerator, making $Q < K$ and forcing a shift toward the reactants.

  \end{core}

  \begin{core}[label=core:letemp]{Le Chatelier's Principle (\textit{Temperature Change})}
    
    A temperature shift is analogous to an increase or decrease of heat within a system, and the Q shift direction depends on whether heat is a net product (\textbf{exothermic}) or a reactant (\textbf{endothermic}) of the chemical reaction at hand. Reconsider the following gaseous reaction in the context of heat.
    
    \bce{\ce{2NO2(g) <=> N2O4(g) + heat}}

    Much like the addition of \ce{N2O4}, increasing the temperature drives the equilibrium to the left to consume the excess energy. Conversely, cooling the system shifts the reaction to the right as it attempts to regenerate the lost heat.

    Chemical equilibrium in the context of temperature is eleaborated more upon in \autoref{sec:shifting_eq}.

  \end{core}

\section{Types of Equilibria}

  Equilibria can also be categorized into types.

  \begin{core}{Homogeneous Equilibria}
    
    In \textbf{homogeneous equilibria}, all reactants and products are in the same phase (solid, liquid, or gas), such as the following examples.

    \bce{\ce{Fe^{3+}(aq) + SCN^{-}(aq) <=> FeSCN^{2+}(aq)}\hspace{5mm}Solution}

    \bce{\ce{H2(g) + I2(g) <=> 2HI(g)}\hspace{5mm}Gas}

    \bce{\ce{HA(aq) <=> H^{+}(aq) + A^{-}(aq)}\hspace{5mm}Acid-Base}

  \end{core}
  
  In \textbf{heterogeneous equilibria}, the reactants and products are in different phases and can be further subcategorized into three domains.  

  \begin{core}{Heterogeneous Equilibria (\textit{Solubility of Solids})}
    
    The solubility of solids equilibrium refers to the dissociation and crystallization rate of solutes. Consider the dissolution of glucose. Solutes that do not dissociate are called \textbf{molecular or nonelectrolyte solutes}.

    \bce{\ce{C6H12O6(s) <=> C6H12O6(aq)}}

    \bce{\text{Dissociation Rate} = $k_f(SA)$}

    \bce{\text{Crystallization Rate} = $k_c(SA)$[\ce{C6H12O6(aq)}]}

    The dissociation rate only depends on the surface area (SA) of the solid glucose, while the crystallization rate depends on both the SA and [\ce{C6H12O6}]. At equilibrium:
    \[
      K_{sp} = \frac{k_c(SA)}{k_f(SA)} = [\ce{C6H12O6(aq)}]
    \]

    For glucose, the $K_{sp}$ value is 5.1M at 25°C, meaning that the maximum [\ce{C6H12O6}] in water is 5.1M. Beyond this concentration, any excess glucose will remain undissolved.

    Now consider a different compound like \ce{NaCl}. Solutes that dissociate are called \textbf{ionic or eletrolyte solutes}.

    \bce{\ce{NaCl(s) <=> Na^{+}(aq) + Cl^{-}(aq)}}

    \bce{\text{Dissociation Rate} = $k_f(SA)$}

    \bce{\text{Crystallization Rate} = $k_c(SA)$[\ce{Na^{+}(aq)}][\ce{Cl^{-}(aq)}]}

    At equilibrium:
    \[
      K_{sp} = \frac{k_c(SA)}{k_f(SA)} = [\ce{Na^{+}(aq)}][\ce{Cl^{-}(aq)}]
    \]

    At 25$^\circ$C, the $K_{sp} \approx 37.3$, making the maximum [\ce{NaCl}] in water approximately $6.11M$.

  \end{core}
  
  \begin{core}{Heterogeneous Equilibria (\textit{Solubility of Gases}) \& Henry's Law}

    Gases can disolve into liquids just as how solids can. Consider the $K_{sp}$ of \ce{CO2} in water.

    \bce{\ce{CO2(g) <=> CO2(aq)}}

    When we reach equilibrium, we have the following:
    \[
      K_{H} = \frac{[\ce{CO2(aq)}]}{\ce{P_{CO2}}} \to [\ce{CO2(g)}] = K_{H} \ce{P_{CO2}}
    \]

    The equilibrium formula above is known as \textbf{Henry's Law}, and the constant $K_H$ is called the \textbf{Henry's Constant} ($M/atm$), a quantitative measure of the proportion between a gas's solubility in a liquid and its partial pressure. The higher the $K_H$ value, the easier it is for a gas to dissolve in a liquid.

    \myfig{0.8}{equilibrium_figures/gas_eq.png}{Gas solubility dynamic equilibrium}{fig:gaseq}

  \end{core}  

  \begin{core}{Heterogeneous Equilibria (\textit{Phase Change})}

    Phase changes also experience equilibriums. Consider the vaporization of water in a sealed flask. 

    \bce{\ce{H2O(l) <=> H2O(g)}}

    \bce{\text{Evaporation Rate} = $k_f(SA)$}

    \bce{\text{Condensation Rate} = $k_c(SA)$ \ce{P_{\ce{H2O}}}}

    Just as the solubility equilibria for nonelectrolyte solutes is only dependent on the solvent concentration, the phase change equilibria for a gas to liquid transition is only dependent on \ce{P_{\ce{H2O}}}.
    \[
      K_{sp} = \frac{k_c(SA)}{k_f(SA)} = \ce{P_{\ce{H2O}}}
    \]

    \myfig{0.7}{equilibrium_figures/phase_eq.png}{Phase change equilibrium of liquid-gas systems}{fig:phaseeq}

    \textit{Phase change equilibria become irrelevant for solid-liquid transitions because both pure solids and pure liquids have fixed, constant concentrations that are excluded from the equilibrium expression, leaving no variable quantities and rendering the equilibrium constant meaningless.}

  \end{core}

\section{Shifting Equilibria}
\label{sec:shifting_eq}

  In \textbf{phase change equilibria}, the melting and boiling points of \ce{H2O} at $0^\circ$C and $100^\circ$C respectively are defined at a standard pressure of 1 atm; under varying pressure conditions, these transition temperatures become dynamic.

  \myfig{0.5}{equilibrium_figures/water_phases.png}{Phase diagram of water}{fig:waterphases}

  Temperature shifts have a more nuanced effect in the context of equilibrium compared to the simple Le Chatelier's principle we have seen in \autoref{core:letemp}. Namely, the $K_{sp}$ values of various solutes change across temperature ranges.


  On the molecular scale, $K_{sp}$ values of solid solutes increase with higher temperatures because the increased molecular motion of the solute and solvent molecules raises the probability of the solute entering the aqueous state.

  In the context of gas solubility within a closed system, the Henry's Law constant $K_{H}$ decreases with increasing temperature. Since gaseous molecules possess greater kinetic energy than their aqueous counterparts, the dissolution of a gas is an exothermic process; as temperature rises, the equilibrium shifts to favour the gaseous state, increasing partial pressure $P$ at the expense of aqueous concentration $C$, thereby reducing $K_{H} = \frac{C}{P}$.
  
  \myfig{0.4}{equilibrium_figures/s_curve.png}{Solubility curve of various compounds}{fig:scurve}

\end{document}

